\subsection{Fourierreihe}
\begin{emphbox}
    Die komplexe Fourierreihe eines periodischen Signals $x(t)$ lautet
    \begin{equation*}
        \begin{array}{@{}cc@{}}
            \text{Synthese}                                   & \text{Analyse}                                     \\
            x(t)=\sum_{k=-\infty}^{\infty}c_ke^{jk\omega_0 t} & c_k=\frac{1}{T_0}\int_0^{T_0}x(t)e^{-jk\omega_0 t}
        \end{array}
    \end{equation*}
\end{emphbox}

\begin{sectionbox}
    \subsubsection{Umrechnungen zu $\cos$/$\sin$-Darstellung}
    $a_k=2\Re{c_k}=c_k+c_{-k} \quad b_k=-2\Im{c_k}=c_k-c_{-k}$ \\
    $c_k=\frac{1}{2}(a_k-jb_k) \qquad c_{-k} = c_k^*=\frac{1}{2}(a_k+jb_k)$
\end{sectionbox}

\begin{sectionbox}
    \subsubsection{Eigenschaften}
    \begin{itemize}
        \item \textbf{Halbwellensymmetrie 1. Art}: Falls $a_0=0$ und $x(t)=-x(t-\frac{T}{2})$, so gilt $a_{2k}=b_{2k}=0$
        \item $\frac{a_0}{2}$ entspricht dem konstanten Gleichanteil, um die Fourierreihe zum Signaloffset auf der y-Achse zu verschieben
    \end{itemize}
\end{sectionbox}

\begin{sectionbox}
    \subsubsection{Parsevalsche Theorem}
    Leistungsäquivalenz zwischen Leitsungsberechnung im Zeitbereich und Frequenzbereich
    \begin{equation*}
        P=\frac{1}{RT}\int_0^Tx^2(t)\diff t = \frac{1}{R}\sum_{k=-\infty}^{\infty}|c_k|^2=
    \end{equation*}
    \begin{equation*}
        =\frac{1}{R}\left(\left(\frac{a_0}{2}\right)^2 + \frac{1}{2}\sum_{k=1}^\infty(a_k^2+b_k^2)\right)
    \end{equation*}
    mit $\frac{1}{R} |c_k|^2$: Leistung der $k$-ten Schwingung. Mit $2$ multiplizieren, falls $|c_k|=|c_{-k}|$
\end{sectionbox}